% chapters/16/rangkuman.tex
\bagian{Rangkuman}

\begin{rangkumanbox}
\begin{itemize}
  \item Pembuat permainan tidak dapat menilai karyanya sendiri dengan baik,
        sebab ia sudah tahu tombol mana yang harus ditekan dan jalur mana yang
        harus dilewati. Pengetahuan itu justru menutupi masalah.
  \item Ada dua jenis pengujian yang saling melengkapi. Pemeriksaan otomatis
        menemukan hal yang bersifat pasti dan cepat; \istilah{playtest}
        menemukan hal yang tidak dapat dihitung.
  \item Pemeriksaan otomatis dijalankan lebih dahulu, sebab tidak ada gunanya
        meminta orang lain menguji level yang belum dapat diselesaikan. Namun
        ia tidak dapat menggantikan \istilah{playtest}.
  \item Pengujian adalah siklus, bukan langkah terakhir. Setiap perbaikan
        menuntut pemeriksaan ulang, sebab perbaikan dapat menimbulkan temuan
        baru.
  \item \istilah{Checklist} fungsional diturunkan dari GDD, dan butirnya
        sebaiknya terbagi menjadi tiga kelompok: mekanik, integrasi, dan
        keseluruhan. Kelompok ketiga hanya dapat diperiksa pada
        \istilah{build}.
  \item \istilah{Soft-lock} adalah keadaan ketika pemain tidak dapat
        melanjutkan dan tidak dapat kembali, walaupun programnya masih
        berjalan dan tidak ada pesan kesalahan.
  \item Jangkauan lompatan dapat dihitung dari besaran fisika, bukan ditebak.
        Dengan $v_0 = 700$ dan $g = 1800$ piksel, tinggi lompatannya 4,25
        petak dan jarak mendatar terjauhnya 7,29 petak pada ketinggian yang
        sama.
  \item Algoritma pemeriksaan level sama dengan algoritma pencarian jalur pada
        Bab~X. Yang berbeda hanya aturan tetangganya: dari empat petak
        bersebelahan menjadi pijakan mana pun yang dapat dicapai dengan
        berjalan, jatuh, atau melompat.
  \item Pemeriksaan keterjangkauan dan simulasi dapat memberi jawaban yang
        sangat berbeda. Pada level contoh, keduanya mengembalikan 15,6 persen
        dan 82,0 persen. Angka yang tinggi menyesatkan, sebab ia hanya
        menyentuh keadaan yang kebetulan terpilih.
  \item Pemeriksaan otomatis tidak dapat menilai apakah levelnya menyenangkan,
        tidak memeriksa bunyi maupun animasi, dan hanya memeriksa level yang
        tertulis pada berkasnya.
  \item Lembar \istilah{playtest} memuat waktu dan tempat, bukan hanya kesan.
        Waktu yang tercatat memudahkan Anda menyempitkan dugaan atas temuan
        yang tidak dapat diulang.
  \item Catatan \istilah{bug} yang berguna memuat tiga hal: apa yang terjadi,
        apa yang seharusnya terjadi, dan urutan tindakan yang menyebabkannya.
  \item Temuan ditimbang atas dua pertanyaan: apakah permainannya masih dapat
        diselesaikan, dan berapa banyak pemain yang akan menemukannya.
  \item Temuan yang berulang pada beberapa sesi biasanya menandakan GDD yang
        belum jelas, bukan kesalahan pelaksanaan.
  \item Bab ini menuntaskan \SUBCPMK{4.3}, dan dengannya seluruh materi
        mata kuliah ini selesai. Yang tersisa adalah menilai kembali seluruh
        capaian Anda pada Bab~XVII.
\end{itemize}
\end{rangkumanbox}
